\section{阈值为什么会移动：数据、后训练与度量轴}

若把 emergence 理解成固定参数量触发器，就无法解释为什么不同模型家族在相同任务上跨阈值的时间不同。本章沿课件顺序拆解三种移动机制：更好的数据可让能力更早出现；针对性 fine-tuning 可把已发现行为压到更小模型；换用 loss 或 perplexity（困惑度，即 held-out cross-entropy 交叉熵的指数）作为横轴，则可能把“不同大小”重新解释为“不同有效质量”。

\subsection{能力前沿不是一条固定竖线}

课件用手绘能力前沿把任务空间想象成随模型规模扩大的区域。图中的虚线只是某个观察位置，深色区域表示较小模型做不到而较大模型做得到的任务，白色区域则表示当前模型仍未解决的任务；它不是宇宙规定的 100B boundary。这个概念图要回答的是“能力集合怎样扩张”，不是为每项任务宣布一个不可移动的参数门槛，因此必须与后面的 data、fine-tuning 和 model-quality 证据一起阅读。

\lecturefigure{slide-13.jpg}{语言模型规模与可达任务区域的手绘框架}{官方课件第 13 页；Wei et al., 2022}

\begin{knowledgebox}{读图：颜色与白区的课堂解释}
深蓝区域代表小模型已经能完成的任务，新增色带代表跨过当前规模后才可达的 emergent abilities；右上白区不是“理论上不可解”，而是尚未被当前模型解决。学生追问白区是否必须依赖 100B 参数，讲者明确回答：阈值只是已有模型上的观察，更好的 data、architecture 与 algorithm 完全可能向左移动边界。
\end{knowledgebox}

\begin{importantbox}{发现能力与压缩能力是两个问题}
大模型常被用作 capability discovery engine：它揭示哪些行为在足够通用的系统中可能存在。一旦行为被明确描述，研究者可以收集或生成针对性数据，再用 fine-tuning、distillation 或工具化方法让更小模型学会它。前者面对 unknown unknowns，后者面对已知目标。
\end{importantbox}

\subsection{Better data：相关性证据与受控证据}

课堂首先比较 LaMDA、GPT-3 与 PaLM，观察到 PaLM 在更小参数量就出现某项能力。讲者的 running hypothesis 是 PaLM 数据更好，但他立即承认大型预训练成本使模型间差异无法被完全消融。因此，模型家族比较提供线索，不能单独证明数据是唯一原因。

\lecturefigure{slide-14.jpg}{更好的训练数据可让能力在更小模型上出现}{官方课件第 14 页；Wei et al., 2022}

\begin{warningbox}{不能把跨模型家族比较写成数据因果证明}
PaLM 与 LaMDA/GPT-3 同时在 data mixture、architecture、tokenizer、optimization 和 compute allocation 上不同。讲者把数据质量称为主要假设，但没有声称完成受控 ablation。高质量读图必须同时写出 observed shift 与 causal uncertainty。
\end{warningbox}

为了得到更强的因果证据，下一张图转向小型可控实验：固定模型规模和大部分训练语料，只改变某些 verb 在训练集中的出现频率，然后观察 subject--verb agreement error。此时横轴更接近真正被操纵的变量。

\lecturefigure{slide-15.jpg}{受控实验：in-domain verb frequency 降低语法错误}{官方课件第 15 页；Wei et al., 2022}

\begin{knowledgebox}{读图：大模型相关性与小模型消融互补}
横轴是目标 verb 在训练数据中的频率，纵轴是 error rate；随着 in-domain exposure 增加，错误下降。这个 toy result 不足以解释 PaLM 的全部 emergence，却证明“在 compute、model size 和其余数据近似固定时，任务相关数据能移动能力曲线”。它为大型模型比较提供了机制上可信的支撑。
\end{knowledgebox}

\teachervoice{课堂问答中，有学生追问 PaLM 的差异究竟来自数据还是架构。讲者没有虚构一个干净答案，而是说明大型预训练无法经济地完成全套消融；随后用较小、可控的 verb-frequency 实验展示能被证明的那一部分。}

\subsection{Fine-tuning：把已知 desired behavior 压到小模型}

当研究者已经知道目标是 instruction following、step-by-step reasoning 或某种风格时，可以直接优化这个行为。课件借助 RLHF / instruction-following 结果说明，小模型经针对性后训练后可以在目标分布上优于更大的弱后训练模型。

\lecturefigure{slide-16.jpg}{针对 desired behavior 的 fine-tuning 可让小模型超过更大 base model}{官方课件第 16 页；InstructGPT 相关结果}

\begin{knowledgebox}{读图：模型大小不是唯一排序键}
先比较同一种 training recipe 下的模型规模，再比较同一模型规模下的 post-training 方法。绿色标记的小型 RLHF/instruction-tuned model 在人类偏好指标上超过更大 base model，说明目标行为可以通过高信息密度监督被提前诱导；但它不代表小模型在所有未知任务上获得了大模型的完整能力集合。
\end{knowledgebox}

学生问能否把 emergent ability distill 到小模型，讲者的回答是肯定的：大模型可以生成数据，小模型再在这些数据上 fine-tune。限制在于研究者必须先知道要生成什么行为；未被发现的能力无法被逐项列入训练目标。

\subsection{横轴应该是参数、FLOPs，还是 perplexity}

参数量和训练 FLOPs 便于公开比较，却只是 model quality 的代理。若数据与算法显著改善，一个更小模型可能有更低 held-out loss，并在 downstream task 上更接近较大旧模型。因此，课堂提出把 perplexity 或 dev-set loss 放到横轴，观察任务能力是否更一致地随有效预测质量变化。

\lecturefigure{slide-17.jpg}{不同 scale 轴：参数、训练 FLOPs、loss 与 perplexity}{官方课件第 17 页；Wei et al., 2022}

\begin{knowledgebox}{背景概念：cross-entropy 与 perplexity}
令真实 token 为 $x_t$，模型给出的条件概率为 $p_\theta(x_t\mid x_{<t})$。平均 cross-entropy 为
\[
H=-\frac{1}{T}\sum_{t=1}^{T}\log p_\theta(x_t\mid x_{<t}),
\]
perplexity 定义为
\[
\mathrm{PPL}=\exp(H).
\]
直觉上，PPL 可理解为模型在每一步面对的“有效等概率选项数”；越低通常越好。但它受 tokenizer、evaluation corpus 与 domain mismatch 影响，不能跨任意设置无条件比较，也不能替代任务正确率。
\end{knowledgebox}

\begin{warningbox}{Perplexity 不是理解力总分}
在 WikiText-103 上更会预测下一个 token，不保证在医学、代码或多步推理上更可靠。把 PPL 放到横轴的价值，是减少“同参数量但训练质量不同”的混淆；它仍是特定 held-out distribution 上的平均 proxy。
\end{warningbox}

\teachervoice{讲者在课堂上把 perplexity 解释为 held-out 文本上的 next-word prediction 能力，并承认未来若小模型用更好数据和算法达到更低 perplexity，按参数量和按模型质量画出的 emergence 图可能明显不同。}

\subsection{Few-shot 超越 fine-tuning 的社会学变化}

当一个通用 pretrained model 只靠上下文中的几个 exemplar，就能超过某些为单任务训练的 state-of-the-art system，研究工作流也随之改变：团队不必先为每个任务收集数千 labels，而可以先测试统一模型的 interface。这是课件称为 ``surprising finetuning'' 的含义。

\lecturefigure{slide-18.jpg}{规模扩大后，few-shot prompting 在若干任务上超过先前 fine-tuned baseline}{官方课件第 18 页；Wei et al., 2022}

\begin{knowledgebox}{读图：跨越绿色 baseline 意味着什么}
绿色线是先前 task-specific fine-tuned result，蓝色点是随规模变化的 few-shot pretrained model。跨越只证明在该 benchmark、该 prompt 和该成本定义下，通用模型的 accuracy 超过旧 baseline；它不证明 fine-tuning 失去价值，因为生产系统还要考虑 latency、serving cost、privacy、稳定性和大量 in-domain labels。
\end{knowledgebox}

课堂把这一点与产品工程作了清晰区分。面向未知任务的 general model 追求一套权重覆盖广泛能力，允许昂贵计算；Google Translate 这类已知任务拥有大量平行数据和严格 latency 约束，往往更适合小型专用模型。``更通用'' 与 ``更适合部署'' 不是同一排序。

\subsection{本章小结}

能力阈值可以被 better data、targeted fine-tuning、distillation、architecture 和 metric choice 移动。参数量只是最醒目的坐标，不是唯一原因。研究上最重要的区分是：跨模型家族的相关性提供现象，小规模受控 ablation 提供因果线索，生产约束决定最终选择。

